\chapter{Výskumná štúdia}
\label{chap:vyskum}

Návrh systému predpokladá, že jednotlivé dimenzie prezentačných schopností majú rôzny vplyv na vnímané hodnotenie prezentácie a že sústredené výskyty negatívnych javov sú pre publikum relevantnejšie ako rovnomerne roztrúsené. Aby som tieto predpoklady empiricky overil a nastavil váhy hodnotenia na základe skutočného ľudského vnímania, realizoval som používateľskú štúdiu so skupinou respondentov.

\section{Výskumné otázky a hypotézy}

Štúdia sa zameriava na dve výskumné otázky:

\textbf{RQ1:} Majú vizuálne dimenzie (očný kontakt, pohyb tela) a hlasové dimenzie (hlas) väčší vplyv na vnímané celkové hodnotenie prezentácie ako plynulosť reči (výplňové slová)? Platí hierarchia vizuálne~>~hlasové~>~plynulosť?

\textbf{RQ2:} Je sústredený výskyt negatívneho javu vnímaný horšie ako rovnaký celkový čas javu rozloženého rovnomerne počas celej prezentácie?

Na základe týchto otázok boli definované štyri hypotézy:

\begin{itemize}
    \item \textbf{H1.1a} --- Vizuálne a hlasové dimenzie majú spolu väčší vplyv na celkové hodnotenie ako plynulosť reči.
    \item \textbf{H1.1b} --- Platí hierarchia vplyvu: vizuálne~>~hlasové~>~plynulosť.
    \item \textbf{H1.2} --- Monotónny hlas je vnímaný horšie ako rovnako intenzívne výplňové slová.
    \item \textbf{H2.1} --- Sústredený výskyt výplňových slov (v jednom bloku) je vnímaný horšie ako rovnomerne roztrúsený.
\end{itemize}

\section{Dizajn experimentu}

\subsection{Príprava videí}

Pre štúdiu bolo nahrané séria šiestich videí s rovnakým prezentujúcim a rovnakým obsahom. Každé video malo dĺžku približne 60 sekúnd. Videá sa líšili práve tým aspektom prezentácie, ktorý bol predmetom skúmania. Manipulácia vždy zasiahla iba jednu dimenziu, ostatné boli ponechané neutrálne.

\begin{table}[H]
    \centering
    \begin{tabularx}{\textwidth}{|l|l|X|}
        \hline
        \textbf{Video} & \textbf{Označenie} & \textbf{Popis manipulácie} \\
        \hline
        \specialrule{1.2pt}{0pt}{0pt}
        V1 & Baseline & Neutrálna prezentácia bez zámerných chýb. Slúži ako referenčná hodnota. \\
        \hline
        V2 & Zlý očný kontakt & Prezentujúci sa pozerá na poznámky a mimo kamery --- simulácia absencie očného kontaktu. \\
        \hline
        V3 & Pohyb bokov & Výrazné kývanie bokov zo strany na stranu počas celej prezentácie. \\
        \hline
        V4 & Monotónny hlas & Prezentujúci hovorí s konštantnou výškou a rytmom hlasu, bez akejkoľvek prozódie. \\
        \hline
        V5 & Výplňové slová (rovnomerne) & Výplňové slová sú rozložené rovnomerne počas celej prezentácie (9 výskytov, $\approx$8 za minútu). \\
        \hline
        V6 & Výplňové slová (sústredene) & Rovnaký celkový počet výplňových slov sústredený do strednej časti prezentácie (6 výskytov, $\approx$5,7 za minútu). \\
        \hline
    \end{tabularx}
    \caption{Prehľad videí použitých v štúdii}
    \label{tab:videos}
\end{table}

\subsection{Zber dát}

Dáta boli zozbierané prostredníctvom online dotazníka (Google Forms). Respondentom bolo postupne zobrazených všetkých šesť videí, pričom poradie videí bolo rovnaké pre všetkých. Po každom videu hodnotili prezentujúceho na škále 0--10 v piatich dimenziách: \textbf{celkové hodnotenie}, \textbf{očný kontakt}, \textbf{pohyb tela}, \textbf{plynulosť reči} a \textbf{hlas}. Celkový počet respondentov bol \textbf{n~=~52}.

\section{Štatistické metódy}

Pre overenie hypotéz boli použité nasledujúce štatistické metódy, implementované v jazyku Python s využitím knižnice \texttt{pingouin}\cite{pingouin}:

\begin{itemize}
    \item \textbf{Viacnásobná lineárna regresia} --- štandardizované beta koeficienty vyjadrujúce relatívny vplyv každej dimenzie na celkové hodnotenie. Dáta boli pred regresiou štandardizované (z-score).
    \item \textbf{Pearsonova korelácia per respondent} --- pre každého respondenta bola vypočítaná korelácia medzi hodnoteniami dimenzií a celkovým hodnotením naprieč šiestimi videami. Takto vzniklo 52 párov hodnôt pre každú dimenziu.
    \item \textbf{Párový t-test} --- porovnanie priemerov dvoch skupín merané na tých istých respondentoch\cite{rmAnova}. Použitý pre H1.1, H1.2 a H2.1.
    \item \textbf{Wilcoxonov test} --- neparametrická alternatíva k párovému t-testu\cite{fieldStatistics}. Slúži ako robustnostná kontrola pri porušení normality.
    \item \textbf{Cohenovo d}\cite{cohenD} --- veľkosť efektu (malý: $|d| < 0{,}2$, stredný: $|d| \approx 0{,}5$, veľký: $|d| > 0{,}8$).
    \item \textbf{Repeated Measures ANOVA} s Greenhouse-Geisserovou korekciou\cite{greenhouseGeisser} --- testuje, či existujú štatisticky významné rozdiely medzi všetkými šiestimi videami súčasne. Post-hoc párové testy s Bonferroniho korekciou identifikujú konkrétne rozdielne dvojice.
    \item \textbf{Spearmanova rank korelácia}\cite{spearmanRank} --- porovnanie poradia videí podľa modelu s poradím podľa ľudských hodnotení.
\end{itemize}

Pred spustením RM ANOVA boli overené jej predpoklady: normalita pomocou Shapiro-Wilkovho testu\cite{shapiroWilk}, sféricita pomocou Mauchlyho testu\cite{mauchlySphericity} a prítomnosť odľahlých hodnôt podľa IQR pravidla. Normalita bola mierne porušená pre V1 ($p < 0{,}001$) a V5 ($p = 0{,}047$), čo je pri $n = 52$ akceptovateľné vzhľadom na Centrálnu limitnú vetu\cite{fieldStatistics}. Sféricita bola porušená ($W = 0{,}476$, $p < 0{,}001$, $\varepsilon = 0{,}780$), preto bola použitá Greenhouse-Geisserova korekcia.

\section{Výsledky}

\subsection{H1.1 --- Vplyv dimenzií na celkové hodnotenie}

Štandardizované beta koeficienty z viacnásobnej lineárnej regresie:

\begin{table}[H]
    \centering
    \begin{tabular}{|l|c|}
        \hline
        \textbf{Dimenzia} & \textbf{Beta koeficient} \\
        \hline
        \specialrule{1.2pt}{0pt}{0pt}
        Hlas (voice) & 0,843 \\
        \hline
        Očný kontakt (eye) & 0,586 \\
        \hline
        Plynulosť (fluency) & 0,574 \\
        \hline
        Pohyb tela (body) & 0,250 \\
        \hline
    \end{tabular}
    \caption{Štandardizované beta koeficienty z lineárnej regresie}
    \label{tab:betas}
\end{table}

Empirická hierarchia je teda \textbf{hlas~>~očný kontakt~>~plynulosť~>~pohyb tela}, čo je v rozpore s predpokladanou hierarchiou vizuálne~>~hlasové~>~plynulosť.

\textbf{H1.1a:} Párový t-test na per-respondent koreláciách (Fisherovou z-transformáciou): $t = -2{,}143$, $p = 0{,}037$. Priemerná korelácia kombinovanej skupiny vizuálne+hlasové ($\bar r = 0{,}550$) je štatisticky signifikantne \textbf{nižšia} než korelácia plynulosti s celkovým hodnotením ($\bar r = 0{,}613$). \textit{Hypotéza bola zamietnutá v opačnom smere} --- vizuálne+hlasové dimenzie spolu nie sú silnejším prediktorom než plynulosť. Príčinou je nízky príspevok pohybu tela ($\beta = 0{,}250$), ktorý ``stiahne'' priemer kombinovanej skupiny pod úroveň plynulosti.

\textbf{H1.1b:} Vizuálne vs. hlasové: $t = -5{,}238$, $p < 0{,}001$ --- hlasové dimenzie \textbf{štatisticky signifikantne dominujú} nad vizuálnymi. Hlasové vs. plynulosť: $t = 1{,}457$, $p = 0{,}151$ --- rozdiel nie je signifikantný. \textit{Hypotéza bola čiastočne podporená v opačnom smere:} hlas dominuje nad vizuálnymi dimenziami, ale poradie hlas~>~plynulosť nebolo štatisticky dokázané.

\textbf{Metodologická poznámka --- halo efekt:} Pri within-subjects dizajne hodnotili respondenti každé video vo všetkých dimenziách súčasne. V dátach sa prejavil halo efekt: keď bolo video zámerne zhoršené v jednej dimenzii, klesali aj hodnotenia ostatných dimenzií, ktoré neboli manipulované (napr. V4 s monotónnym hlasom mal nižšie hodnotenie plynulosti $5{,}67$ oproti baseline $8{,}18$). Beta koeficienty z regresie teda odrážajú aj tento efekt, nie len izolovaný vplyv každej dimenzie, a ich interpretáciu ako čistých váh treba brať s rezervou.

\subsection{H1.2 --- Monotónnosť vs. výplňové slová}

Priemerné celkové hodnotenie V4 (monotónny hlas) = 4,40~±~1,91 a V5 (výplňové slová rovnomerne) = 4,44~±~2,09. Párový t-test: $t = -0{,}164$, $p = 0{,}870$, $d = -0{,}019$. Wilcoxon: $W = 386{,}0$, $p = 0{,}742$.

\textit{Hypotéza nebola podporená.} Oba javy boli respondentmi hodnotené identicky. Pravdepodobným vysvetlením je intenzita manipulácie v V5 --- výplňové slová boli natoľko frekventované, že spôsobili rovnako silný negatívny dojem ako monotónny hlas.

\subsection{H2.1 --- Distribúcia výplňových slov}

Priemerné celkové hodnotenie V5 (rovnomerne) = 4,44~±~2,09 a V6 (sústredene) = 6,27~±~2,07. Párový t-test: $t = 5{,}492$, $p < 0{,}001$, $d = 0{,}871$. Wilcoxon: $W = 81{,}0$, $p < 0{,}001$. 40 z 52 respondentov (76,9\,\%) hodnotilo V5 horšie ako V6.

\textit{Hypotéza bola zamietnutá s opačným nálezom.} Rovnomerne rozložené výplňové slová sú vnímané výrazne horšie ako sústredené. Konštantná prítomnosť výplňových slov neposkytuje poslucháčovi žiadnu ``oddychovú fázu'', čo vedie k vyššej miere rušenia. Efekt je veľký ($d = 0{,}871$) a robustný.

\subsection{Repeated Measures ANOVA}

RM ANOVA potvrdila, že medzi videami existujú štatisticky vysoko významné rozdiely: $F(5, 255) = 39{,}412$, $p_{GG} < 10^{-22}$, $\eta^2_g = 0{,}256$ (veľký efekt). Bonferroniho post-hoc testy odhalili 12 z 15 signifikantných párov ($p < 0{,}05$). Nesignifikantné páry (V4 vs. V5, V2 vs. V6, V3 vs. V6) sú v súlade s ostatnými nálezmi.

RM ANOVA per dimenzia odhalila najväčší efekt pre \textbf{hlas} ($\eta^2_g = 0{,}402$), nasledovaný plynulosťou ($\eta^2_g = 0{,}334$) a očným kontaktom ($\eta^2_g = 0{,}175$). Pohyb tela mal najmenší efekt ($\eta^2_g = 0{,}050$). Tieto hodnoty sú v súlade s beta koeficientmi z regresie.

\section{Implikácie pre systém}

Výsledky štúdie boli zapracované do systému v dvoch krokoch:

\textbf{Model v2 --- aktualizácia váh:} Pôvodné váhy dimenzií (hlas 25\,\%, plynulosť 20\,\%, pohyb tela 25\,\%, očný kontakt 30\,\%) boli nahradené váhami odvodenými z normalizovaných beta koeficientov (súčet beta = 2,24):

\begin{table}[H]
    \centering
    \begin{tabular}{|l|c|c|}
        \hline
        \textbf{Dimenzia} & \textbf{Model v1} & \textbf{Model v2} \\
        \hline
        \specialrule{1.2pt}{0pt}{0pt}
        Hlas & 25\,\% & 37\,\% \\
        \hline
        Plynulosť & 20\,\% & 27\,\% \\
        \hline
        Očný kontakt & 30\,\% & 27\,\% \\
        \hline
        Pohyb tela & 25\,\% & 9\,\% \\
        \hline
    \end{tabular}
    \caption{Porovnanie váh dimenzií medzi modelmi}
    \label{tab:weights}
\end{table}

\textbf{Model v3 --- penalizácia distribúcie výplňových slov:} Na základe nálezu H2.1 bol do výpočtu skóre plynulosti pridaný distribučný trest. Prezentácia je rozdelená do 6 časových okien; koeficient variácie (CV = $\sigma / \mu$) hustoty výplňových slov určuje mieru rovnomernosti distribúcie. Nízke CV (rovnomerné rozloženie) vedie k vyššej penalizácii (max. 20 bodov).

\subsection{Porovnanie modelov}

Validita modelov bola overená Spearmanovou rank koreláciou medzi poradím videí podľa systému a poradím podľa ľudských hodnotení ($n = 52$):

\begin{table}[H]
    \centering
    \begin{tabular}{|l|c|c|}
        \hline
        \textbf{Model} & \textbf{Spearman $\rho$} & \textbf{Zmena} \\
        \hline
        \specialrule{1.2pt}{0pt}{0pt}
        v1 (pôvodné váhy) & 0,771 & --- \\
        \hline
        v2 (váhy z výskumu) & 0,943 & $+0{,}171$ \\
        \hline
        v3 (váhy + distribúcia) & 0,943 & $0{,}000$ \\
        \hline
    \end{tabular}
    \caption{Spearmanová rank korelácia modelov s ľudskými hodnoteniami}
    \label{tab:spearman}
\end{table}

Model v2 dosiahol výrazné zlepšenie oproti modelu v1 --- poradie videí zodpovedá ľudskému vnímaniu výrazne presnejšie. Model v3 nezmenil celkové poradie (V5 bolo pred V6 už v modeli v2), avšak zväčšil rozdiel skóre V5 a V6 v dimenzii plynulosť z 16,2 na 22,4 bodu. Analýza ukázala, že systémový rozdiel v modeli v3 ($\Delta = -2{,}24$) presahuje ľudský rozdiel ($\Delta = -1{,}64$) --- distribučná penalizácia je teda mierne agresívnejšia, ako zodpovedá ľudskému vnímaniu, a jej kalibrácia ostáva predmetom ďalšieho vývoja.

Systémové skóre naprieč všetkými modelmi systematicky nadhodnocuje kvalitu prezentácie v porovnaní s ľudskými hodnoteniami (priemerný rozdiel $\approx +27$ bodov). Príčinou je, že systém deteguje len merateľné technické javy, zatiaľ čo ľudia vnímajú prezentáciu holisticky --- vrátane aspektov, ktoré systém neanalyzuje (nervozita, dikcia, obsah). Kalibrácia absolútnych hodnôt skóre zostáva otvorenou výzvou pre ďalší vývoj.

\section{Validácia hodnotiaceho modelu}

Tretia výskumná oblasť overuje, či poradie videí podľa systémového skóre zodpovedá poradiu podľa ľudského hodnotenia a či empiricky odvodené váhy a distribučná penalizácia túto zhodu zlepšujú. Analýza je realizovaná v \texttt{validation.ipynb} a opiera sa o Spearmanovu rank koreláciu medzi 6 systémovými skóre a 6 ľudskými priemermi.

\subsection{H3.1 --- Empiricky odvodené váhy zvyšujú zhodu modelu s ľudským hodnotením}

Hypotéza: Spearmanova korelácia poradia videí bude vyššia pre model v2 (empirické váhy) ako pre model v1 (intuitívne váhy).

Výsledky sú zhrnuté v tabuľke~\ref{tab:spearman} v predchádzajúcej sekcii. Prechod z modelu v1 ($\rho = 0{,}771$, $p = 0{,}072$) na model v2 ($\rho = 0{,}943$, $p = 0{,}005$) predstavuje nárast o $+0{,}171$. Model v1 nie je štatisticky signifikantný, model v2 áno. \textit{Hypotéza H3.1 je potvrdená.}

Ako doplnková kontrola bola vykonaná leave-one-out krížová validácia (LOO-CV): pre každé z 6 videí boli váhy odvodené z regresie na zostávajúcich 5 videách a aplikované na vynechané video. Výsledná korelácia dosiahla $\rho = 0{,}829$ ($p = 0{,}042$, $\text{MAE} = 0{,}43$), čo naznačuje, že regresne odvodené váhy generalizujú aj mimo trénovacej vzorky.

\subsection{H3.2 --- Distribučná penalizácia presnejšie kopíruje ľudský rozdiel V5 vs. V6}

Hypotéza: Absolútna chyba $|\Delta_\text{sys} - \Delta_\text{human}|$ na dimenzii plynulosť bude nižšia pre model v3 ako pre modely v1 a v2.

Ľudský rozdiel hodnotenia plynulosti V5~$-$~V6: $\Delta_\text{human} = -1{,}64$ (párový t-test: $t = -5{,}864$, $p < 0{,}001$).

\begin{table}[H]
    \centering
    \begin{tabular}{|l|c|c|}
        \hline
        \textbf{Model} & \textbf{Systémový $\Delta$ (V5\,$-$\,V6)} & \textbf{$|\Delta_\text{sys} - \Delta_\text{human}|$} \\
        \hline
        \specialrule{1.2pt}{0pt}{0pt}
        v1 & $-1{,}62$ & 0,015 \\
        \hline
        v2 & $-1{,}62$ & 0,015 \\
        \hline
        v3 & $-2{,}24$ & 0,605 \\
        \hline
    \end{tabular}
    \caption{Porovnanie systémového a ľudského rozdielu V5~$-$~V6 na dimenzii plynulosť}
    \label{tab:h32}
\end{table}

\textit{Hypotéza H3.2 nebola potvrdená --- výsledok je opačný.} Modely v1 a v2 sú výrazne bližšie k ľudskému rozdielu ($|\text{chyba}| = 0{,}015$) ako model v3 ($|\text{chyba}| = 0{,}605$). Distribučná penalizácia zveličuje rozdiel medzi rovnomerne a sústrednene rozloženými výplňovými slovami nad mieru, akú ľudia skutočne vnímajú. Smer efektu (V5 vnímané horšie ako V6) je zachytený správne vo všetkých troch modeloch, avšak veľkosť penalizácie v modeli v3 si vyžaduje ďalšiu kalibráciu.

\subsection{Obmedzenia validácie}

Validácia je realizovaná na rovnakom korpuse 6 videí, ktoré boli použité aj v štúdii RQ1/RQ2. Pri $n = 6$ videách je Spearmanova korelácia citlivá na jednotlivé body a bootstrap intervaly spoľahlivosti pri tejto veľkosti vzorky neposkytujú konfirmačnú informáciu. LOO-CV zostáva tiež v rovnakom korpuse. Validácia na nezávislom korpuse reálnych prezentácií je predmetom ďalšieho výskumu.
